% !TeX program = pdflatex
% Corrige TD 2 - Analyse L1 - 2MTH1 0101
\documentclass[11pt,a4paper,oneside]{article}
\usepackage[T1]{fontenc}
\usepackage[utf8]{inputenc}
\usepackage[main=french,provide=*]{babel}
\usepackage{lmodern,microtype,amsmath,amssymb,mathtools}
\usepackage{geometry}
\geometry{left=2.4cm,right=2.3cm,top=2.55cm,bottom=2.45cm,headheight=24pt}
\usepackage{enumitem,booktabs,needspace}
\usepackage[most]{tcolorbox}
\usepackage{fancyhdr,lastpage,titlesec,hyperref}
\definecolor{Marine}{HTML}{10283F}
\definecolor{Bleu}{HTML}{176D8E}
\definecolor{Sarcelle}{HTML}{118F86}
\definecolor{Or}{HTML}{EDAA47}
\definecolor{GrisTexte}{HTML}{54667A}
\definecolor{BleuPale}{HTML}{EFF5F8}
\hypersetup{colorlinks=true,linkcolor=Bleu,urlcolor=Sarcelle,pdftitle={Analyse L1 - Solutions TD 2 - Calcul differentiel},pdfauthor={M. Bakary TRAORE}}
\setlength{\parindent}{0pt}
\setlength{\parskip}{5pt plus1pt minus1pt}
\setlist[enumerate]{leftmargin=1.55em,itemsep=3pt,topsep=3pt}
\titleformat{\section}{\Large\bfseries\color{Marine}}{}{0pt}{}[\vspace{1mm}{\color{Sarcelle}\titlerule[.8pt]}]
\titleformat{\subsection}{\large\bfseries\color{Bleu}}{}{0pt}{}
\titlespacing*{\section}{0pt}{19pt}{8pt}
\titlespacing*{\subsection}{0pt}{12pt}{4pt}
\pagestyle{fancy}\fancyhf{}
\fancyhead[L]{\small\color{Marine}\textbf{UNIVERSITÉ DE N'ZÉRÉKORÉ}}
\fancyhead[R]{\small\color{GrisTexte}Analyse L1 \textbullet\ 2MTH1 0101}
\fancyfoot[L]{\scriptsize\color{GrisTexte}Corrigé TD 2 \textbullet\ Calcul différentiel}
\fancyfoot[R]{\scriptsize\color{GrisTexte}\thepage\,/\,\pageref{LastPage}}
\newtcolorbox{objectifbox}{enhanced,breakable,colback=BleuPale,colframe=Bleu,boxrule=.6pt,arc=1mm}
\newcommand{\R}{\mathbb R}
\newcommand{\e}{\mathrm e}
\newcommand{\solution}{\textbf{\color{Sarcelle}Solution.}\enspace}
\begin{document}
\begin{tcolorbox}[enhanced,colback=Bleu,colframe=Bleu,boxrule=0pt,sharp corners,left=5mm,right=5mm,top=5mm,bottom=5mm]
\begin{center}
{\color{white}\large\bfseries UNIVERSITÉ DE N'ZÉRÉKORÉ\par}\vspace{.14cm}
{\color{white!85}\normalsize LICENCE 1 -- TRONC COMMUN \quad / \quad 2MTH1 0101\par}\vspace{.30cm}
{\color{white}\LARGE\bfseries SOLUTIONS DÉTAILLÉES DU TD 2\par}\vspace{.14cm}
{\color{white!95}\large CHAPITRE 2 -- Calcul différentiel\par}\vspace{.14cm}
{\color{white!80}\small M. Bakary TRAORE \quad | \quad 2026--2027\par}
\end{center}
\end{tcolorbox}
\begin{objectifbox}
\textbf{Usage.} Les dix exercices de la feuille TD 2 sont repris dans leur ordre, avec justifications, calculs, conditions d'application et interprétations. À lire après un essai autonome.
\end{objectifbox}

\section*{Partie A -- Dérivation et dérivées successives}
\Needspace{6\baselineskip}
\subsection*{Exercice 1 -- Règles de dérivation}
\textit{Énoncé :} déterminer domaines et dérivées de $(x^2+1)\e^{-x}$, $\ln(1+x^2)$ et $\sqrt{2x+3}/(x-1)$.

\solution
\begin{enumerate}[label=\alph*)]
\item $f$ est définie sur $\R$. Produit et chaîne :
\[f'(x)=2x\e^{-x}-(x^2+1)\e^{-x}=\boxed{-(x-1)^2\e^{-x}}.\]
\item $1+x^2>0$ pour tout $x\in\R$. Donc $D_g=\R$ et
\[\boxed{g'(x)=\frac{2x}{1+x^2}}.\]
\item $D_h=[-\frac32,+\infty[\,\setminus\{1\}$. Pour $x>-\frac32$, $x\ne1$,
\begin{align*}
h'(x)&=\frac{(x-1)(1/\sqrt{2x+3})-\sqrt{2x+3}}{(x-1)^2}\\
&=\boxed{\frac{-x-4}{(x-1)^2\sqrt{2x+3}}}.
\end{align*}
La formule n'est pas définie au bord $x=-\frac32$ ; ce point n'est pas intérieur au domaine.
\end{enumerate}

\Needspace{6\baselineskip}
\subsection*{Exercice 2 -- Dérivée par définition et tangente}
\textit{Énoncé :} $f(x)=x^2$ en $x=2$ ; $u(x)=|x|$ en $0$.

\solution Pour $h\ne0$,
\[\frac{f(2+h)-f(2)}h=\frac{4+4h+h^2-4}{h}=4+h\longrightarrow4.\]
Ainsi $f'(2)=4$, $f(2)=4$ et l'équation de la tangente est
\[\boxed{y=4+4(x-2)=4x-4}.\]
Pour $u(x)=|x|$, le quotient $|h|/h$ vaut $1$ si $h>0$, et $-1$ si $h<0$. Les limites unilatérales sont différentes : $u$ n'est pas dérivable en $0$.

\Needspace{6\baselineskip}
\subsection*{Exercice 3 -- Dérivées successives}
\textit{Énoncé :} dériver $f(x)=\e^{2x}$ et démontrer la formule générale.

\solution On trouve $f'=2\e^{2x}$, $f''=4\e^{2x}$ et $f'''=8\e^{2x}$. L'égalité
\[\boxed{f^{(n)}(x)=2^n\e^{2x}}\]
est vraie au rang $0$. Si elle est vraie au rang $n$, alors
\[(f^{(n)})'=2^n\cdot2\e^{2x}=2^{n+1}\e^{2x}.\]
Elle est donc vraie pour tout $n\in\mathbb N$ par récurrence.

\section*{Partie B -- Théorèmes et extrema}
\Needspace{6\baselineskip}
\subsection*{Exercice 4 -- Théorème de Rolle}
\textit{Énoncé :} $f(x)=x^2-4x+3$ sur $[1,3]$.

\solution Un polynôme est continu sur $[1,3]$ et dérivable sur $]1,3[$. On vérifie
\[f(1)=1-4+3=0,\qquad f(3)=9-12+3=0.\]
Par Rolle, il existe $c\in]1,3[$ tel que $f'(c)=0$. Or $f'(x)=2x-4$, donc $\boxed{c=2}$.

\Needspace{6\baselineskip}
\subsection*{Exercice 5 -- Théorème des accroissements finis}
\textit{Énoncé :} $f(x)=\ln x$ sur $[1,\e]$.

\solution Le logarithme est continu sur $[1,\e]$ et dérivable sur $]1,\e[$. Sa pente moyenne vaut
\[\frac{\ln\e-\ln1}{\e-1}=\frac1{\e-1}.\]
Le TAF assure un $c\in]1,\e[$ tel que $1/c=1/(\e-1)$, soit
\[\boxed{c=\e-1\approx1.7183}\quad\text{et}\quad1<c<\e.\]
La pente moyenne est celle de la droite sécante joignant les extrémités du graphe.

\Needspace{6\baselineskip}
\subsection*{Exercice 6 -- Étude complète d'une fonction}
\textit{Énoncé :} $f(x)=x^3-3x$ sur $[-2,2]$.

\solution 
\begin{enumerate}[label=\alph*)]
\item $f'(x)=3(x+1)(x-1)$ : le signe est $+$ sur $]-\infty,-1[$, $-$ sur $]-1,1[$, $+$ sur $]1,+\infty[$. Donc maximum local $\boxed{f(-1)=2}$ et minimum local $\boxed{f(1)=-2}$.
\item Aux bornes, $f(-2)=-2$ et $f(2)=2$. Les \textbf{maxima globaux} sur $[-2,2]$ valent $2$, atteints en $x=-1$ et $x=2$ ; les \textbf{minima globaux} valent $-2$, atteints en $x=-2$ et $x=1$.
\item $f''(x)=6x$. Elle est négative sur $]-2,0[$ et positive sur $]0,2[$ : $f$ est concave sur $[-2,0]$, convexe sur $[0,2]$, avec \textbf{point d'inflexion} $\boxed{(0,0)}$.
\end{enumerate}

\section*{Partie C -- DL et approximations}
\Needspace{6\baselineskip}
\subsection*{Exercice 7 -- Développements limités}
\textit{Énoncé :} deux limites, approximation de $\e^{-0.1}$ et erreur.

\solution
\begin{enumerate}[label=\alph*)]
\item On utilise les DL $\e^x=1+x+\frac{x^2}2+o(x^2)$ et $\ln(1+x)=x-\frac{x^2}2+o(x^2)$. Après division par $x^2$,
\[\boxed{\lim_{x\to0}\frac{\e^x-1-x}{x^2}=\frac12},\qquad
\boxed{\lim_{x\to0}\frac{\ln(1+x)-x}{x^2}=-\frac12}.\]
\item À l'ordre 2 : $\e^{-0.1}\approx 1-0.1+\frac{0.1^2}{2}=\boxed{0.905}$. Par Taylor--Lagrange avec $M=\e^{0.1}$,
\[\boxed{|R_2|\le\frac{\e^{0.1}(0.1)^3}{6}<0.000185}.\]
En réalité $\e^{-0.1}\approx0.9048374$ et l'erreur $\approx0.0001626$ respecte la borne.
\end{enumerate}

\section*{Partie D -- Applications aux quatre filières}
\Needspace{6\baselineskip}
\subsection*{Exercice 8 -- Météorologie et hydrologie}
\textit{Énoncé :} étudier $T(t)=25+3\sin(\pi t/12)$ et $Q(t)=5+2t\e^{-t}$.

\solution
\begin{enumerate}[label=\alph*)]
\item $T'(t)=\frac\pi4\cos(\pi t/12)$. Donc $\boxed{T'(0)=\pi/4\approx0.7854\ ^\circ\mathrm C/\mathrm h}$ : au début de l'observation, la température augmente à cette vitesse instantanée.
\item $Q'(t)=2(1-t)\e^{-t}$. Comme $\e^{-t}>0$, la dérivée est positive sur $[0,1[$, nulle en $1$, négative pour $t>1$. Le maximum a lieu à $\boxed{t=1\ \mathrm h}$ et vaut $\boxed{Q(1)=5+2/\e\approx5.736\ \mathrm{m^3/s}}$.
\end{enumerate}

\Needspace{6\baselineskip}
\subsection*{Exercice 9 -- Génie de l'environnement}
\textit{Énoncé :} $C(t)=10\e^{-0.2t}$ en mg/L.

\solution $C'(t)=10(-0.2)\e^{-0.2t}=\boxed{-2\e^{-0.2t}=-0.2C(t)}$. Ainsi $C'(0)=-2\ \mathrm{(mg/L)/h}$. À $t=0.5$ h, $-0.2t=-0.1$ :
\[C(0.5)\approx10\left(1-0.1+\frac{0.1^2}{2}\right)=\boxed{9.05\ \mathrm{mg/L}}.\]
La concentration est décroissante pour $t\ge0$ puisque $C'<0$ ; il s'agit d'une diminution dans ce modèle idéalisé.

\Needspace{6\baselineskip}
\subsection*{Exercice 10 -- Gestion des Ressources Naturelles}
\textit{Énoncé :} $B(t)=100-20\e^{-0.1t}$ en tonnes.

\solution Pour $t\ge0$,
\[\boxed{B'(t)=2\e^{-0.1t}>0},\qquad
\boxed{B''(t)=-0.2\e^{-0.1t}<0}.\]
La biomasse augmente, mais sa vitesse de croissance diminue. $B(0)=80$ tonnes et $\lim_{t\to+\infty}B(t)=\boxed{100\ \mathrm{tonnes}}$ : le stock tend vers un niveau asymptotique dans ce modèle.

\begin{objectifbox}
\textbf{Méthode générale.} Vérifier le domaine, appliquer une règle de calcul adaptée, exploiter le signe de la dérivée, écrire les unités, puis interpréter le résultat sans confondre le modèle mathématique et une observation réelle.
\end{objectifbox}
\end{document}
